\documentclass[a4paper,12pt]{article}
\usepackage[T2A]{fontenc}
\usepackage[utf8]{inputenc}
\usepackage[russian]{babel}
\usepackage[margin=24mm]{geometry}
\usepackage{amsmath,amssymb,booktabs,tabularx,xcolor,tikz,fancyhdr}
\definecolor{accent}{HTML}{008080}
\definecolor{darkaccent}{HTML}{004D4D}
\usepackage[colorlinks=true,linkcolor=darkaccent]{hyperref}
\setlength{\parindent}{0pt}
\setlength{\parskip}{6pt}
\setlength{\headheight}{16pt}
\pagestyle{fancy}
\fancyhf{}
\fancyhead[L]{\small Степень с натуральным показателем}
\fancyhead[R]{\small Кирилл Зиновьев}
\fancyfoot[C]{\thepage}
\newcommand{\err}[1]{\textcolor{red!75!black}{\textbf{Ошибка:}} #1}
\newcommand{\idea}[1]{\textcolor{blue!60!black}{\textbf{Ключевая идея:}} #1}
\newcommand{\result}[1]{\textcolor{darkaccent}{\textbf{Итог:}} #1}
\newcommand{\parthead}[1]{\par\medskip\textbf{#1}\par}
\begin{document}
\hypersetup{pageanchor=false}
\begin{titlepage}
\centering
\vspace*{25mm}
{\Large\color{darkaccent} МАТЕМАТИКА\par}
\vspace{14mm}
{\huge Ревью домашней работы\par}
\vspace{9mm}
{\LARGE\color{accent} Степень с натуральным показателем\par}
\vfill
{\large Кирилл Зиновьев\par}
\vspace{8mm}
27 ноября 2026 г.\par
\vspace{14mm}
{\small Лёвин Артём Александрович\\эксклюзивно для Кирилла Зиновьева\par}
\vspace{17mm}
\begin{tikzpicture}
\draw[accent,line width=.7pt] (0,0) .. controls (2,.65) and (4,-.65) .. (6,0) .. controls (8,.65) and (10,-.65) .. (12,0);
\end{tikzpicture}
\end{titlepage}
\hypersetup{pageanchor=true}
\section*{\color{darkaccent} Сводная таблица}
Нажмите на номер задания для перехода к подробному разбору.

\medskip
{\renewcommand{\arraystretch}{1.5}
\begin{tabularx}{\linewidth}{@{}c >{\raggedright\arraybackslash}p{3cm} >{\raggedright\arraybackslash}X >{\raggedright\arraybackslash}p{2cm} >{\raggedright\arraybackslash}p{2.6cm}@{}}
\toprule
№ & Тема & Суть ошибки & Ваш ответ & Правильный ответ\\
\midrule
\hyperref[task:2]{2} & Знак и скобки & Потерян минус перед четвёртой степенью & $337$ & $-175$\\
\hyperref[task:8]{8} & Куб десятичной дроби & Неверное основание; куб не записан & $-0{,}003$ & $(-0{,}3)^3$\\
\bottomrule
\end{tabularx}}
\clearpage
\thispagestyle{empty}
\null\vfill
\begin{center}{\Huge\color{darkaccent} Разбор ошибок}\end{center}
\vfill\null
\clearpage
\phantomsection\label{task:2}
\section*{\color{darkaccent} Задание № 2. Знак перед степенью}
\parthead{Текст задания}
Вычислите $-4^4+(-3)^4$.
\parthead{Ваше решение и ответ}
\[-4^4+(-3)^4=256+81=337.\]
Ваш ответ: $337$.

\err{Потеря отрицательного знака перед первой степенью.}

\parthead{Где именно возникает ошибка}
Исходное выражение переписано верно. Первый неверный переход: $-4^4=256$. Вторая степень вычислена верно: $(-3)^4=81$. Сложение полученных положительных чисел также выполнено верно, но первое слагаемое имеет другой знак.

\parthead{Почему этот переход неверен}
Возведение в степень выполняется раньше применения внешнего минуса. В записи $-4^4$ основание степени равно $4$, поэтому минус сохраняется перед результатом. Скобки включают знак в основание: в записи $(-3)^4$ основание равно $-3$. Четыре отрицательных множителя дают положительное произведение.
\[
-4^4=-(4\cdot4\cdot4\cdot4)=-256,\qquad
(-4)^4=256.
\]
\idea{Перед вычислением определите границы основания каждой степени.}

\parthead{Исправление и полное решение}
\begin{align*}
4^4&=16\cdot16=256, & (-3)^4&=9\cdot9=81,\\
-4^4+(-3)^4&=-256+81=-(256-81)=-175.
\end{align*}
\parthead{Проверка}
$256>81$, поэтому сумма отрицательного и положительного слагаемых отрицательна. Разность модулей равна $175$; также $175+81=256$.

\result{$\boxed{-175}$. Внешний минус сохраняйте после вычисления степени.}
\parthead{Задача на закрепление}
Вычислите $-5^4+(-2)^4$.
\clearpage
\phantomsection\label{task:8}
\section*{\color{darkaccent} Задание № 8. Куб десятичной дроби}
\parthead{Текст задания}
Представьте $-0{,}027$ в виде куба десятичной дроби.
\parthead{Ваше решение и ответ}
Записано $-0{,}003$. Промежуточные вычисления и показатель степени отсутствуют.

\err{Выбрано неверное основание, итоговая запись куба отсутствует.}

\parthead{Где именно возникает ошибка}
Отрицательный знак сохранён правильно. Первый достоверно неверный результат --- число $-0{,}003$ в качестве ответа. По записи невозможно определить, каким способом оно получено.

\parthead{Почему этот результат неверен}
Куб числа --- произведение трёх одинаковых множителей. Для десятичной дроби при умножении количество знаков после запятой складывается. Например, у числа $0{,}003$ три десятичных знака; в его кубе их девять:
\[(-0{,}003)^3=-\frac{27}{10^9}=-0{,}000000027.\]
Это значение отличается от заданного числа.

\idea{Разложите числитель и знаменатель десятичной дроби на кубы.}
\parthead{Исправление и полное решение}
Перепишем исходное число как обыкновенную дробь. Числитель $27=3^3$, знаменатель $1000=10^3$. Куб отрицательного числа отрицателен, поэтому
\[
-0{,}027=-\frac{27}{1000}
=\frac{(-3)^3}{10^3}
=\left(-\frac{3}{10}\right)^3
=(-0{,}3)^3.
\]
\parthead{Проверка}
\[(-0{,}3)\cdot(-0{,}3)\cdot(-0{,}3)=0{,}09\cdot(-0{,}3)=-0{,}027.\]
Получено исходное число. У трёх множителей по одному десятичному знаку, поэтому в произведении три знака после запятой.

\result{$\boxed{(-0{,}3)^3}$. Записывайте ответ в требуемом виде степени.}
\parthead{Задача на закрепление}
Представьте $-0{,}064$ в виде куба десятичной дроби.
\end{document}
